\begin{figure}[!htbp]
  \centering
  \begin{tikzpicture}[x=1.18cm,y=4.85cm,>=Stealth]
    % axes
    \draw[->, line width=1.0pt, draw=black!80] (0,0) -- (6.35,0)
      node[below, font=\small] {$k$};
    \draw[->, line width=1.0pt, draw=black!80] (0,0) -- (0,1.06)
      node[left, font=\small] {$\Gamma(k)$};

    % ticks
    \foreach \y/\lab in {0.0/0,0.4/0.4,0.8/0.8,1.0/1.0} {
      \draw[black!75] (-0.04,\y) -- (0.04,\y);
      \node[left, font=\scriptsize, text=black!72] at (-0.08,\y) {\lab};
    }
    \foreach \x in {1,2,3,4,5,6} {
      \draw[black!75] (\x,0.0) -- (\x,-0.018);
      \node[below, font=\scriptsize, text=black!72] at (\x,-0.02) {\x};
    }

    % threshold lines
    \draw[densely dashed, gray!35] (0,0.80) -- (6.0,0.80);
    \draw[densely dashed, gray!55] (0,0.85) -- (6.0,0.85);
    \draw[densely dashed, gray!35] (0,0.90) -- (6.0,0.90);
    \node[right, font=\scriptsize, text=gray!70!black] at (6.03,0.80) {$\rho=0.80$};
    \node[right, font=\scriptsize, text=gray!78!black] at (6.03,0.85) {$\rho=0.85$};
    \node[right, font=\scriptsize, text=gray!70!black] at (6.03,0.90) {$\rho=0.90$};

    % filled areas
    \fill[cyan!10]
      (0.55,0.00) --
      plot[smooth] coordinates {
        (0.55,0.24) (1.00,0.50) (1.55,0.76) (2.05,0.86)
        (2.85,0.93) (3.90,0.97) (5.60,1.00)
      } --
      (5.60,0.00) -- cycle;

    \fill[orange!10]
      (0.55,0.00) --
      plot[smooth] coordinates {
        (0.55,0.10) (1.05,0.22) (1.75,0.37) (2.55,0.52)
        (3.35,0.67) (4.15,0.79) (4.85,0.86) (5.55,0.92)
      } --
      (5.55,0.00) -- cycle;

    % curves
    \draw[cyan!65!black, line width=1.7pt]
      plot[smooth] coordinates {
        (0.55,0.24) (1.00,0.50) (1.55,0.76) (2.05,0.86)
        (2.85,0.93) (3.90,0.97) (5.60,1.00)
      };

    \draw[orange!82!black, line width=1.7pt]
      plot[smooth] coordinates {
        (0.55,0.10) (1.05,0.22) (1.75,0.37) (2.55,0.52)
        (3.35,0.67) (4.15,0.79) (4.85,0.86) (5.55,0.92)
      };

    % highlighted target and k* markers
    \fill[cyan!70!black] (2.05,0.85) circle (1.6pt);
    \fill[orange!82!black] (4.85,0.85) circle (1.6pt);
    \draw[densely dotted, cyan!65!black, line width=0.9pt] (2.05,0.00) -- (2.05,0.85);
    \draw[densely dotted, orange!82!black, line width=0.9pt] (4.85,0.00) -- (4.85,0.85);
    \fill[cyan!65!black] (2.05,0.00) circle (1.35pt);
    \fill[orange!82!black] (4.85,0.00) circle (1.35pt);

    % labels
    \node[font=\scriptsize, text=cyan!70!black,
          fill=white, inner sep=2pt, rounded corners=3pt]
      at (1.95,0.97) {集中型画像};
    \node[font=\scriptsize, text=orange!82!black,
          fill=white, inner sep=2pt, rounded corners=3pt]
      at (5.02,0.60) {弥散型画像};

    \node[font=\scriptsize, text=gray!75!black,
          fill=white, inner sep=2pt]
      at (5.08,0.875) {相同 coverage 目标};

    \node[font=\scriptsize, text=cyan!70!black] at (2.05,-0.08) {$k^\star_{\text{集中}}$};
    \node[font=\scriptsize, text=orange!82!black] at (4.85,-0.08) {$k^\star_{\text{弥散}}$};
  \end{tikzpicture}
  \caption{敏感度覆盖度曲线示意：集中型画像与弥散型画像的预算含义。相同覆盖阈值下，集中型画像会在更小的 $k$ 上达到目标覆盖度，而弥散型画像需要更大的保护层数；这正是 AutoK 将 profile 形状转写为预算建议的核心直觉。}
  \label{fig:ch3-coverage-curve}
\end{figure}
